\section{Frozen source ledger and typed comparison}
\label{sec:source-ledger}

Fix $q\in\{2,3,5\}$ and let
\[
  \Sigma_q=\F_q^{\mathbb Z}
\]
with the left shift $\sigma$. A nonempty word is \emph{primitive} when it is
not a proper power. The source primitive object is its oriented cyclic class:
cyclic rotations are identified, but reversal is not identified unless it is
already a rotation. We write $\Prim_q$ for these primitive necklaces and
$n(\gamma)$ for the length of $\gamma\in\Prim_q$. Ordinary word powers encode
temporal repetition.

\begin{definition}[Frozen source fields]
For $\gamma\in\Prim_q$ of length $n$, define
\[
 T_q(\gamma)=n\log q,
 \qquad
 F_\gamma(s,z)=\left(1-z^nq^{-ns}\right)^{-1}.
\]
The free variable $z$ marks one original shift symbol. Thus the $r$th
repetition contributes $z^{nr}q^{-nrs}$; marker degree and temporal
repetition are distinct pieces of data.
\end{definition}

The weighted adjacency of the one-vertex, $q$-loop graph acts on $\C$ by
\[
  \Lop f=zq^{1-s}f.
\]
It is a scalar operator, but it still owns a complete dynamical determinant.
Indeed, the full shift has exactly $q^r$ points fixed by $\sigma^r$. The
marked periodic-point exponential is therefore
\begin{align}
 Z_q(s,z)
 &=\exp\left(\sum_{r\ge1}\frac{q^r(zq^{-s})^r}{r}\right) \\
 &=\exp\left(\sum_{r\ge1}\frac{(zq^{1-s})^r}{r}\right)
 =\frac{1}{1-zq^{1-s}}.                                      \label{eq:source-zeta}
\end{align}
This is a formal identity in $z$ and an analytic identity in the disk where
the displayed logarithmic series converges. Its inverse is the ordinary
finite-dimensional determinant
\begin{equation}
  \Dq(s,z)=\det(I-\Lop)=1-zq^{1-s}.                           \label{eq:source-det}
\end{equation}
The periodic-point and determinant organization follows the standard finite-
shift framework \citep{artin1965periodic,bowen1970zeta}.

Let $\Nq(n)$ denote the number of primitive cyclic classes of length $n$.
Every word fixed by $\sigma^r$ repeats a unique primitive necklace of length
$d\mid r$, and each such necklace has $d$ based representatives. Hence
\begin{equation}
 q^r=\sum_{d\mid r}d\Nq(d),
 \qquad
 \Nq(n)=\frac1n\sum_{d\mid n}\mu(d)q^{n/d}.                 \label{eq:necklace}
\end{equation}
Regrouping repetitions in \cref{eq:source-zeta} gives
\begin{equation}
 Z_q(s,z)=\prod_{\gamma\in\Prim_q}
 \left(1-z^{n(\gamma)}q^{-s n(\gamma)}\right)^{-1}.          \label{eq:source-product}
\end{equation}
The count in \cref{eq:necklace} also counts monic irreducible polynomials of
degree $n$ over $\F_q$ \citep{chebolu2026necklace}. Under the finite-field
norm, both types carry $q^n$ and the degree clock $n\log q$. This is the
source positive control. We use only count equality; no canonical objectwise
bijection is asserted.

The rational-prime comparator has a different primitive type. Let $\Primes$
be the positive rational primes and define, on $\ell^2(\Primes)$,
\[
 Q_se_p=p^{-s}e_p.
\]
For $\RePart s>1$, the operator is trace class because
$\sum_p|p^{-s}|<\infty$. It owns
\begin{equation}
 \DP(s,z)=\det(I-zQ_s)
 =\prod_{p\in\Primes}(1-zp^{-s}),                             \label{eq:target-det}
\end{equation}
whose reciprocal specializes at $z=1$ to the rational Euler product
\citep{nistDLMF25211}. For sufficiently small $|z|$,
\begin{equation}
 -\log\DP(s,z)
 =\sum_{r\ge1}\frac{z^r}{r}P(rs),
 \qquad
 P(s)=\sum_{p\in\Primes}p^{-s}.                              \label{eq:target-log}
\end{equation}
The diagonal comparator proves that the target product has a valid operator;
it does not make that operator source-owned.

Any credited factorwise descent must preserve several necessary fields at
once. Comparing one source factor with one target factor produces the matrix
in \cref{tab:necessary-fields}.

\begin{table}[t]
\centering
\small
\caption{Necessary fields for an in-place identification. These conditions
are not asserted to be sufficient for a completed rational-prime model.}
\label{tab:necessary-fields}
\begin{tabularx}{\textwidth}{@{}
  >{\raggedright\arraybackslash}p{0.20\textwidth}
  >{\raggedright\arraybackslash}p{0.23\textwidth}
  >{\raggedright\arraybackslash}p{0.23\textwidth}Y@{}}
\toprule
Field & Source length $n$ & Rational-prime target & Forced consequence \\
\midrule
Clock/weight & $n\log q$, $q^{-ns}$ & $\log p$, $p^{-s}$ & $p=q^n$ \\
Primitive marker & $z^n$ & $z$ & $n=1$ \\
Multiplicity & $\Nq(n)$ factors & one factor per $p$ & bijective factor count \\
Repetition & $z^{nr}q^{-nrs}$ & $z^rp^{-rs}$ & only after primitive fields agree \\
Ownership & scalar $q$-loop adjacency & diagonal prime inventory & same owner required \\
\bottomrule
\end{tabularx}
\end{table}

The full-ledger claim also requires totality: every source primitive factor
must be accounted for. A partial map can produce a correct local factor, but
it cannot identify the complete source product with the complete target
product. With these types fixed, the three obstructions in the next section
are direct consequences rather than post-hoc choices of a convenient marker
or projection.
